\documentclass[10pt,a4paper]{amsart}
\usepackage[margin=19mm]{geometry}
\usepackage{amsmath,amssymb,mathtools}
\usepackage[scr=rsfs,cal=euler]{mathalfa}
\usepackage{xfrac}
\usepackage[unicode,hidelinks,bookmarks=false]{hyperref}
\newcommand{\ie}{i.e.\ }
\newcommand{\cf}{cf.\ }
\newcommand{\resp}{respectively\ }
\newcommand{\Subsection}{\subsection}
\providecommand{\nobreakdash}{\nobreak\hbox{-}}
\setlength{\emergencystretch}{2em}
\multlinegap0pt
\usepackage{bm}
\usepackage{bbm}
\usepackage{dsfont}
\usepackage{xspace}
\usepackage{cancel}
\usepackage{mathtools}
\usepackage{amsthm}
\makeatletter
\@ifundefined{thm}{\newtheorem{thm}{Theorem}}{}
\makeatother
\title[Semigroup Congruences and Subsemigroups of the Direct Square]{Semigroup Congruences and Subsemigroups of the Direct Square}
\author{Callum Barber, Nik Ru\v{s}kuc}
\date{Source: arXiv:2504.02428v1 (CC BY 4.0)}
\begin{document}
\maketitle

\noindent\emph{Source and licence.} Semigroup Congruences and Subsemigroups of the Direct Square. Version arXiv:2504.02428v1 (CC BY 4.0), distributed under the Creative Commons Attribution 4.0 International licence, as recorded in the arXiv licence metadata for this exact version and shown on the abstract page https://arxiv.org/abs/2504.02428v1. Full rights notice: licenses/arxiv-2504.02428-CC-BY-4.0.txt.\par\smallskip

\noindent\textit{Editorial scope.} Counted unit: the Thm whose source label is \texttt{Baer Levi} in the article (source0.tex lines 458--460, source label \texttt{Baer Levi}), delivered together with the article's own complete original proof of it (source0.tex lines 462--477, one \texttt{proof} environment of 16 lines). No mathematics is written, repaired or re-derived: statement and proof are the article's own text line for line. Required context retained uncounted: none. Every definition, display and label of those lines is retained unchanged. Omitted from this package and not redistributed: 1--457, 461--461, 478--744. Nothing inside a retained excerpt is omitted or altered: every retained body is delivered exactly as the article's lines are written, so no omission receipt of any kind accompanies this package. Citations: the retained excerpts carry no citation command at all, so no bibliography entry of the article is transcribed and no reference is redistributed. Reference adaptation: none was needed and none was made; every reference inside the delivered unit resolves inside it, so no unresolved reference or citation is produced. Printed slips kept literal: no printed word, symbol, hypothesis or number of the counted statement or of its proof was altered. Layout and class: the article's own class is \texttt{amsart} with packages including amsmath, amssymb, hyperref, amsfonts, mathtools, theoremref, amsthm, xcolor, inputenc, array, enumitem, babel, comment, multicol; the wrapper is the lane's pinned \texttt{amsart} editorial wrapper at 10pt on A4, which restates the theorem-like environments the retained text needs and adds the article's own macro definitions that the retained text needs (none), with extra wrapper packages (bm, bbm, dsfont, xspace, cancel, mathtools). The delivered numbering is the unit's own. Assets: no retained span contains a figure environment, a graphics-inclusion command, a PDF-page inclusion, a table or a TikZ picture; no image file is copied into this package, the package declares no asset dependency and no image file is redistributed with it. Licence: version arXiv:2504.02428v1 of this article is distributed under the Creative Commons Attribution 4.0 International licence, as recorded in the arXiv licence metadata for this exact version and shown on the abstract page https://arxiv.org/abs/2504.02428v1. A full rights notice is delivered at licenses/arxiv-2504.02428-CC-BY-4.0.txt. Characters: every retained excerpt is pure ASCII, so the delivered engine, XeLaTeX with its default Latin Modern text font, reports no missing character.
\begin{thm}\label{Baer Levi}
    For any two infinite cardinals $p \geq q$, the generalized Baer--Levi semigroup $B(p,q)$ is not DSC.
\end{thm}

\begin{proof}
    Let $B = B(p,q)$.
    Consider the set 
    \[
    \rho = \bigl\{ (f,g) \in B \times B \colon (X \setminus X f) \cap (X \setminus X g) \neq \emptyset \bigr\}.
    \]
     It is easy to check that $\rho$ is a diagonal subsemigroup of $B \times B$. It is also clear that $\rho$ is symmetric, so we will show $\rho$ is not transitive.
    Partition $X$ into two sets $A$ and $B$ with cardinality $p$. Let $A^\prime$ and $ B^\prime$ be subsets of $A$ and $B$ respectively both with cardinality $q$.
    Let $x \in A^\prime$. The sets $X \setminus A^\prime , X \setminus (B^\prime \cup {x} )$ and $X \setminus B^\prime$ all have cardinality $p$. 
    So there are bijections 
    \[
    f^\prime\colon X \to X \setminus A^\prime ,\quad  g^\prime \colon X \to X \setminus (B^\prime \cup {x})\quad \text{and}\quad h^\prime \colon  X \to X \setminus B^\prime.
    \] 
    Each can be extended to an injection from $X$ to itself; call these functions $f,g$ and $h$ respectively.
    Note that $X\setminus Xf = A^\prime, X\setminus Xg = B^\prime \cup \{ x \}$ and $X\setminus Xh = B^\prime$. So $f,g,h \in B$ and $(f,g),(g,h) \in \rho$, but $(f,h) \notin \rho$. Hence $\rho$ is not transitive.
\end{proof}

\end{document}
